\section{Auditoria de qualidade do dataset TON\_IoT}
\label{sec:dataset}

Antes de treinar qualquer modelo, fizemos uma auditoria da qualidade
estatística dos sete CSVs de telemetria disponibilizados pela UNSW
(\texttt{Train\_Test\_IoT\_*.csv}). A motivação é simples: o paper
Alsaedi 2020 reporta F1 igual a 1{,}00 em
várias bases (notavelmente \texttt{Garage\_Door} em todos os oito
modelos), e qualquer reprodução que não interrogue \emph{por que}
isso acontece corre o risco de incorporar artefatos não-conceituais e
declarar vitória indevida. Encontramos três problemas estruturais:
dois tipos de \emph{leak} e uma divergência crítica de tamanho entre
os CSVs públicos e o que o paper de fato mediu.

A seção está organizada em cinco partes. A primeira inventaria as
\emph{features} por associação com o rótulo e diagnostica o \emph{leak}
categórico determinístico em duas bases. A segunda caracteriza o
\emph{leak} temporal que satura nas sete bases. A terceira documenta a
divergência de tamanho entre os CSVs públicos e o dataset que produziu
as Tabelas 10--13 do paper. A quarta mostra que os dois \emph{leaks}
determinísticos são, na verdade, um artefato dessa divergência. A
quinta consolida o pipeline mínimo de tratamento que a auditoria torna
obrigatório.

\subsection{Inventário de features e leak categórico}
\label{sec:leak_categorico}

A Figura~\ref{fig:cramers} mostra, para cada base, o maior Cramér's V
entre as \emph{features} contra o rótulo binário (\texttt{normal} vs
qualquer ataque), calculado sobre os dados crus. Valores próximos de 1
indicam uma \emph{feature} deterministicamente associada ao rótulo,
sinal forte de \emph{leak}.

\begin{figure}[ht]
  \centering
  \includegraphics[width=0.95\linewidth]{v06_cramers_v_por_base.png}
  \caption{Cramér's V máximo por base (a \emph{feature} mais associada
           ao rótulo), sobre os dados crus. \texttt{temp\_condition}
           (Fridge) e \texttt{sphone\_signal} (Garage\_Door) atingem
           V $=1{,}000$, uma dependência determinística (leak).
           \texttt{GPS\_Tracker} e \texttt{Weather} têm associação forte
           mas não determinística; as demais bases não têm leak
           marginal.}
  \label{fig:cramers}
\end{figure}

Duas \emph{features} categóricas mostraram Cramér's V exatamente
$1{,}000$ contra o rótulo (Tabela~\ref{tab:leak_categorico}):

\begin{table}[ht]
  \centering
  \caption{Leaks categóricos identificados nos CSVs públicos.}
  \label{tab:leak_categorico}
  \small
  \begin{tabular}{lllc}
    \toprule
    \textbf{Base} & \textbf{Feature} & \textbf{Causa} & \textbf{V} \\
    \midrule
    Fridge       & \texttt{temp\_condition} & 6 variantes \emph{whitespace} de \texttt{'high'}/\texttt{'low'} & 1{,}000 \\
    Garage\_Door & \texttt{sphone\_signal}  & \texttt{'0'/'1'} (atk), \texttt{'false\ \ '/'true\ \ '} (normal) & 1{,}000 \\
    \bottomrule
  \end{tabular}
\end{table}

Em \texttt{Fridge.temp\_condition}, as seis variantes
(\texttt{'high'}, \texttt{'high\ '}, \texttt{'high\ \ '} e as análogas
para \texttt{'low'}) diferem apenas por caracteres de \emph{whitespace}
no final da string. Para um humano elas são a mesma palavra; para um
\texttt{LabelEncoder} ingênuo, são seis classes distintas, e o modelo
aprende a separar \texttt{normal}/\texttt{attack} via essa diferença
ortográfica em vez de via o conceito que a \emph{feature} representa:
a variante \emph{sem} espaço aparece $100\%$ em \texttt{attack}, as
variantes \emph{com} espaço $100\%$ em \texttt{normal}. Análogo em
\texttt{Garage\_Door}, em que o encoding mistura strings booleanas
(\texttt{'false\ \ '}/\texttt{'true\ \ '}, sempre \texttt{normal}) com
strings numéricas (\texttt{'0'}/\texttt{'1'}, sempre \texttt{attack}) de
forma sistemática por classe. A Figura~\ref{fig:leak_summary} torna o
\emph{whitespace} visível (glifo \texttt{\textvisiblespace}) para evidenciar essa
separação.

\begin{figure}[ht]
  \centering
  \includegraphics[width=0.92\linewidth]{v06_leak_summary.png}
  \caption{Sumário dos \emph{leaks} identificados no TON\_IoT
           público: categórico (\texttt{temp\_condition},
           \texttt{sphone\_signal}) e temporal (\texttt{date}).
           O \emph{whitespace} é renderizado com o glifo \texttt{\textvisiblespace}:
           a variante sem espaço é $100\%$ \texttt{attack}, as variantes
           com espaço são $100\%$ \texttt{normal}.}
  \label{fig:leak_summary}
\end{figure}

\subsection{Leak temporal universal}

O paper, em sua Seção VI--A1, descarta \texttt{date}, \texttt{time} e
\texttt{timestamp} sob o argumento de que ``\emph{may cause some ML
methods to overfit}''. Nossa análise mostra uma justificativa mais
forte: o campo \texttt{date} é \textbf{leak determinístico em todas as
sete bases}, com 100\% das datas únicas aparecendo em uma única classe.
A Figura~\ref{fig:date_leak} reporta a fração de datas exclusivas por
classe em cada base; em todas, a barra está fixada em 100\%.

\begin{figure}[ht]
  \centering
  \includegraphics[width=0.92\linewidth]{v06_date_leak_per_base.png}
  \caption{Fração de datas únicas que aparecem em apenas uma classe
           (\texttt{normal} ou \texttt{attack}), por base. As sete
           barras saturam em 100\%: o \emph{testbed} foi
           operado em janelas temporais separadas por classe.}
  \label{fig:date_leak}
\end{figure}

A causa é a operação do \emph{testbed}: os autores rodaram cada
campanha de ataque em uma janela de tempo separada da operação normal.
Para qualquer estudo de generalização, isso significa que \texttt{date}
e variantes derivadas (\texttt{hour}, \texttt{dayofweek}) precisam ser
descartadas \emph{antes} do treinamento, comportamento que
encapsulamos na função \texttt{drop\_temporal\_leak}, aplicada em
todos os experimentos.

\subsection{Divergência de tamanho entre CSVs públicos e o paper}
\label{sec:dataset_mismatch}

O achado mais crítico para a reprodução é que \textbf{os CSVs públicos
\texttt{Train\_Test\_IoT\_*.csv} são um subset re-balanceado do
dataset que produziu as Tabelas 10--13 do paper}. Confrontando as
contagens reais dos CSVs com as Figuras 2/3/5 do paper original
(Tabela~\ref{tab:mismatch}):

\begin{table}[ht]
  \centering
  \caption{Divergência de tamanho entre o dataset usado nas
           Tabelas 10--13 do paper Alsaedi 2020 (lido diretamente das
           Figuras 2/3/5, que reportam \texttt{normal} $=35\,000$ por
           dispositivo e $245\,000$ no \emph{combined}) e os CSVs
           públicos \texttt{Train\_Test\_IoT\_*.csv} (UNSW). Os ataques
           foram preservados; a classe \texttt{normal} foi sub-amostrada
           para $43\%$ do volume original.}
  \label{tab:mismatch}
  \small
  \begin{tabular}{lrrrr}
    \toprule
    & \multicolumn{2}{c}{\textbf{per-device}} & \multicolumn{2}{c}{\textbf{combined}} \\
    \cmidrule(lr){2-3}\cmidrule(lr){4-5}
    & \textbf{normal} & \textbf{total} & \textbf{normal} & \textbf{total} \\
    \midrule
    Paper Alsaedi 2020 (Figs.\ 2/3/5) & 35\,000 & 51--60 k & 245\,000 & 401\,119 \\
    CSVs públicos (UNSW) & 15\,000 & 31--40 k & 105\,000 & 261\,119 \\
    \midrule
    Razão (CSVs / paper)             & \textbf{0{,}43} & -- & \textbf{0{,}43} & -- \\
    \bottomrule
  \end{tabular}
\end{table}

Que os ataques foram preservados \emph{na íntegra} é verificável por
aritmética: as fatias da Figura 5 do paper somam $401\,119$ amostras no
\emph{combined} (das quais $245\,000$ são \texttt{normal}), e
$401\,119 - 245\,000 = 156\,119$ coincide exatamente com a contagem de
ataques que medimos nos CSVs públicos. Apenas a classe \texttt{normal}
encolheu.

A consequência é que mesmo o pipeline mais fiel ao paper não pode
reproduzir os números exatos das Tabelas 10--13: a prevalência de
classe é diferente (paper $\approx$ 39\% ataque / 61\% \texttt{normal};
CSVs $\approx$ 60\% ataque / 40\% \texttt{normal}), o que altera o ponto
de operação de modelos que colapsam
para predição majoritária (LR/LDA/NB em bases pobres como
\texttt{Motion\_Light}) e empurra todos os F-score \emph{weighted}
para baixo (a métrica recompensa proporcionalmente quem acerta a
classe majoritária). Modelos robustos a re-amostragem
(Random Forest, CART) tendem a divergir menos; modelos sensíveis a
volume (NB, LSTM) divergem mais. Voltamos a essa limitação na
Seção~\ref{sec:discussao}.

\subsection{Os leaks determinísticos são artefato do subset re-amostrado}
\label{sec:leak_artifact}

Os quatro \emph{datasets} disponibilizados pela UNSW formam uma
hierarquia: \texttt{telemtry\_data} (logs brutos),
\texttt{Processed\_IoT\_dataset} (versão integral processada,
$\sim$3{,}6 M linhas) e \texttt{Train\_Test\_IoT\_*.csv} (subset
pré-dividido para ML, $\sim$261 k linhas). Confrontando a versão pública
\texttt{Train\_Test} com a versão integral \texttt{Processed}, descobre-se
que os leaks \emph{determinísticos} (V $=1{,}000$, exclusividade
temporal $100\%$) são um artefato \textbf{específico do subset
re-amostrado}: na versão integral, as mesmas variantes de
\emph{whitespace} e as mesmas datas existem, mas a associação com o
rótulo é apenas parcial (Tabela~\ref{tab:leak_subset}).

\begin{table}[ht]
  \centering
  \caption{O leak determinístico é uma propriedade do subset público,
           não do dataset integral. A sub-amostragem da classe
           \texttt{normal} (\S{}\ref{sec:dataset_mismatch}) removeu
           justamente as linhas \texttt{normal} de \texttt{'high'}/\texttt{'low'}
           sem espaço, deixando essas variantes $100\%$ em \texttt{attack}.}
  \label{tab:leak_subset}
  \small
  \begin{tabular}{lcc}
    \toprule
    \textbf{Sinal} & \textbf{\texttt{Train\_Test} (público)} & \textbf{\texttt{Processed} (integral)} \\
    \midrule
    \texttt{temp\_condition} (Cramér's V)  & \textbf{1{,}000} & 0{,}569 \\
    \texttt{sphone\_signal} (Cramér's V)   & \textbf{1{,}000} & 0{,}554 \\
    \texttt{date} (exclusividade de classe) & \textbf{100\%} (7/7 bases) & 50--62\% \\
    \bottomrule
  \end{tabular}
\end{table}

A implicação é que a divergência de tamanho
(\S{}\ref{sec:dataset_mismatch}) e o leak categórico
(\S{}\ref{sec:leak_categorico}) não são achados independentes: são o
\emph{mesmo} fenômeno. Sub-amostrar \texttt{normal} de
$\sim$35\,000 para 15\,000 por dispositivo eliminou exatamente as
poucas linhas \texttt{normal} que usavam a variante sem espaço, o que
transformou um sinal parcial (V $\approx 0{,}56$) em um separador
perfeito (V $=1{,}000$). Qualquer reprodução sobre os CSVs públicos
herda esse artefato.

\subsection{Conclusão da auditoria}

A auditoria define o pipeline mínimo necessário para qualquer
comparação metodologicamente honesta no TON\_IoT público:

\[
  \texttt{df} \rightarrow \texttt{drop\_temporal\_leak} \rightarrow
  \texttt{encode} + \texttt{scale} \rightarrow \texttt{modelo}.
\]

O passo opcional \texttt{clean\_df} (que normaliza
\emph{whitespace} e mapeia strings booleanas) é discutido na
Seção~\ref{sec:discussao}: usá-lo destrói o \emph{leak} categórico,
mas afasta a reprodução do paper, que aparentemente operou sobre
\emph{features} já limpas (a diferença está \emph{nos CSVs públicos}, e
não no dataset original do paper).
